\section{Measuring whether a causal assistant has improved}
\label{exec20:validity-scope}

A bank comparing two causal assistants needs to distinguish a better answer
from an answer that is easier for its grader to approve. An assistant may
produce a polished diagram, an executable adjustment formula and a confident
explanation while changing the population, conditioning on a variable measured
after the decision, or overlooking an unmeasured common cause. A test that
rewards the format can improve while the scientific error remains. Conversely,
an assistant that correctly refuses a nonidentified query can look worse under
a rubric that rewards every numerical response. The evaluation must therefore
specify the scientific service, the observation used to assess it and the
conditions under which the observation measures that service.

Lin's validity-guided workflow separates an LLM used as a research tool, an
evaluation target, a human simulator and a cognitive model. Its six stages
connect the intended claim to instrument development, experimental design,
documented execution, dependence-aware analysis and bounded interpretation.
The source is a methodological tutorial, with illustrative supplementary code;
it does not report an original empirical validation of a banking assistant.
For this proposal, its most useful contribution is the discipline of validating
the measurement before interpreting a model comparison. A classifier requires
evidence that it performs its specified task; fluent explanations do not
establish a psychological mechanism, and simulated customer answers require
separate evidence of correspondence to the intended population.
\citep{exec20_validity}

\subsection{The target is correctness under a declared information contract}
\label{exec20:validity-target}

Consider a retention-offer question whose request identifies the eligible
customers, the offer time, the balance horizon and the information available
before the offer. The assistant must preserve these elements when it proposes
a graph and an estimand. A useful benchmark includes several legitimate outputs:
a supported identification formula, an explicit effect set, or a refusal that
names the missing assumption or observation. Correctness is judged against the
stated graph class and information, rather than against a single preferred
diagram. If the scenario does not determine an edge direction, selecting the
examiner's favorite direction is not the target.

Let $A\in\{0,1\}$ denote the assigned assistant workflow, with the base model,
retrieval access and available tools specified for each arm. Let $C=1$ mean that
an answer meets a frozen semantic criterion. Let $G=1$ mean that a grader marks
the answer as passing. These are distinct binary variables. The criterion
could require a correct adjustment set and an explicit warning about an
unidentified query; other criteria, such as factual grounding and valid code,
should also be reported individually. A single binary endpoint is introduced
here only to expose a measurement problem. It cannot capture every consequence
of a wrong scientific answer.

Define $p_a=\Pr(C=1\mid A=a)$, grader sensitivity
$s_a=\Pr(G=1\mid C=1,A=a)$ and grader specificity
$c_a=\Pr(G=0\mid C=0,A=a)$. All probabilities refer to the same declared
scenario and answer population. By conditioning on correctness,
\begin{equation}
 q_a:=\Pr(G=1\mid A=a)
 =s_a p_a+(1-c_a)(1-p_a)
 =(s_a+c_a-1)p_a+1-c_a .
 \label{exec20:validity-measurement}
\end{equation}
This is an authored measurement model, not an empirical estimate from Lin's
tutorial. Random assignment of $A$, consistency of the assigned workflows and
an appropriate no-interference condition can identify a causal effect on the
observed grade for the experimental population. Those conditions alone do not
make $q_1-q_0$ an effect on semantic correctness. The relation between $G$ and
$C$ must also be supported. \citep{exec20_validity_reference}

\subsection{A better grade can conceal an unchanged answer}
\label{exec20:validity-example}

Suppose true correctness equals $p_0=p_1=3/5$. Impose grader properties
$(s_0,c_0)=(4/5,9/10)$ for the first workflow and
$(s_1,c_1)=(19/20,4/5)$ for the second. The more polished output is easier to
recognize when correct but also more likely to be approved when wrong. These
properties are assumptions of the example, not observations about a particular
grader. Equation~\eqref{exec20:validity-measurement} gives
\begin{equation}
 q_0=\frac{13}{25}=0.52,\qquad
 q_1=\frac{13}{20}=0.65,\qquad
 q_1-q_0=\frac{13}{100},\quad p_1-p_0=0.
 \label{exec20:validity-illusion}
\end{equation}
More samples can estimate the observed 13-percentage-point improvement more
precisely while leaving its interpretation wrong. Randomizing the workflow
does not remove this error because the workflow itself changes the relation
between the answer and its measurement.

If the two calibration quantities are known for each arm and
$s_a+c_a\ne1$, inversion gives
\begin{equation}
 p_a=\frac{q_a+c_a-1}{s_a+c_a-1}.
 \label{exec20:validity-inverse}
\end{equation}
It recovers $3/5$ in both arms of the example. Even an invariant grader need
not preserve the size of an improvement. With $s=4/5,c=9/10$, changing true
correctness from $1/2$ to $7/10$ changes observed passing by $7/50$; dividing
by $s+c-1=7/10$ recovers the true change of $1/5$. When $s+c=1$, the grade
contains no information about $p$: with $s=c=1/2$, true correctness $1/10$
and $9/10$ both produce a pass probability of $1/2$.
\citep{exec20_validity_reference}

Calibration is therefore an information requirement. In practice, $s_a$ and
$c_a$ must be estimated using held-out adjudication under the intended arms,
scenario population and rubric. Both correct and incorrect cases need support;
verification only of accepted answers cannot identify both quantities without
additional sampling information. Uncertainty in the three inputs must propagate
through the ratio in~\eqref{exec20:validity-inverse}. A denominator close to
zero makes correction unstable. Known bounds on calibration can instead be
propagated to a set of compatible correctness rates, intersected with $[0,1]$;
an unsupported point correction should not be substituted. Model updates,
subgroup differences and changes in answer style can invalidate transported
calibration. A finite example with known probabilities proves none of those
transport assumptions.

The economic consequence is immediate but remains conditional. Suppose, solely
for planning, that an additional correct result is worth USD~100 and the new
workflow costs USD~8 more per request. Using true correctness, the incremental
expected value is $100(p_1-p_0)-8=-8$ dollars. Substituting grades produces
$100(q_1-q_0)-8=5$ dollars. The sign reverses without any improvement in the
scientific service. Real costs of wrong advice, useful refusal and delayed
decisions may differ sharply; the example supplies neither a bank price nor an
investment recommendation. \citep{exec20_validity_reference}

\subsection{Repeated answers and the population of inference}
\label{exec20:validity-dependence}

Lin distinguishes individual observations, experimental units and units of
analysis, and recommends pairing answers to shared items and respecting grouped
observations. This guidance must be applied to the actual design. A fixed model
version is not automatically a randomly sampled model from a population, and a
prompt condition is not automatically the smallest randomized unit. If a
workflow is randomized across independent scenarios, the scenario can be the
assignment unit. If both workflows are run on each scenario with controlled
order and fresh contexts, the paired difference can be the analysis unit. These
designs target different sources of variation from repeatedly sampling decoder
outputs for one fixed question. \citep{exec20_validity}

Aggregation preserves the dependence that remains between groups. To make this
qualification explicit, consider an authored real-valued score model
\begin{equation}
 Z_{jr}=U+V_j+e_{jr},\qquad
 \overline Z_j=U+V_j+\frac{1}{R}\sum_{r=1}^R e_{jr},
 \label{exec20:validity-shock}
\end{equation}
where all distinct components are independent, have zero mean, and have finite
variances $\sigma_U^2,\sigma_V^2,\sigma_e^2$. Here $j$ indexes groups and $r$
repeated scores; the common $U$ can represent a shared evaluation-environment
shock. This is a separate real-valued example, not the binary measurement model.
For distinct groups $j,k$,
\begin{equation}
 \operatorname{Var}(\overline Z_j)
 =\sigma_U^2+\sigma_V^2+\frac{\sigma_e^2}{R},\qquad
 \operatorname{Cov}(\overline Z_j,\overline Z_k)=\sigma_U^2.
 \label{exec20:validity-averages}
\end{equation}
With variances $1/4,1/2,1$ and $R=20$, group means have variance $4/5$,
covariance $1/4$ and correlation $5/16$. Averaging has reduced decoder noise
but has not made the groups independent. \citep{exec20_validity_reference}

Accordingly, the tutorial's shorthand about aggregation eliminating correlation
cannot justify independent group means without assumptions about the shared
environment and scenario sampling. No fixed cluster-count threshold guarantees
calibration, and balanced group sizes alone do not establish a universal
equivalence between aggregation and every mixed model. For a fixed list of
models, descriptive paired differences over a declared scenario collection may
be the honest result. Generalization to unseen scenarios requires a sampling or
transport argument; generalization to future model versions requires another.

\subsection{A public evaluation design for the proposed service}
\label{exec20:validity-design}

The proposed first evaluation uses generated causal mechanisms with known
identification status, together with public scenario descriptions whose
permitted uses and limitations are documented. Generated mechanisms allow an
exact answer key for graph compatibility and intervention queries. Public
business descriptions test translation and missing-information detection, but
their prose does not supply hidden causal truth. The two sources should retain
separate outcome denominators and support different claims.

During development, analysts can revise prompts, retrieval rules and the
semantic rubric on a designated subset. The final evaluation then freezes those
choices and withholds whole mechanisms and scenario families, including their
paraphrases and graph-isomorphic variants. This extends the tutorial's
development/held-out distinction to leakage risks specific to causal reasoning.
Formatting perturbations, variable renaming and column order are diagnostic
conditions within a scenario family; they are not new independent domains.
\citep{exec20_validity}

Before running the comparison, specify whether the intervention changes only
the prompt, or also the model, retrieval corpus, numerical engine and repair
budget. A paired design can compare two frozen workflows on the same held-out
scenarios while controlling order and preventing shared conversation state.
The prespecified outcomes should include semantic correctness, supported
refusals, unsupported assertions, parsing failures, computation failures and
coverage of requests. Report all assigned requests, with a justified treatment
of missing or invalid outputs, so that a workflow cannot improve its headline
score merely by discarding difficult cases.

The grading calibration sample should span both arms and include independently
adjudicated correct and incorrect answers under a prespecified verification
sampling scheme. Blind adjudicators to workflow identity where feasible, retain
disagreements and report both the original grade and the corrected or bounded
semantic result. If expert judgment is itself fallible, disagreement and
verification cannot simply be renamed truth. Known generated answers support
some calibration tasks; interpretation of an ambiguous business request may
require a set of admissible answers.

Finally, preserve the exact requests, outputs, parser decisions, model and tool
versions, retrieval snapshots, collection times, token counts and retries.
Report grading and adjudication costs alongside execution cost and decision
quality. A deterministic checker for a formal formula, independent expert
adjudication and an LLM grader are substantive alternatives with different
coverage and error properties. This design specifies what a prototype should
demonstrate using public and generated evidence before an expensive bank trial;
it does not report a completed experiment or validated bank economics.

